%=======================================================================
%  Research Design and Methodology
%=======================================================================
\section{Research Design and Methodology}
\label{sec:methodology}

\subsection{Research strategy: why simulation}

The research questions ask what \emph{would} happen to an election
under attack. That rules out the two obvious empirical strategies. A
field experiment is unethical and illegal: it would require attempting
to corrupt a real election. Observational study of historical fraud
suffers from severe selection bias, because only detected fraud is
observable and detected fraud is by definition the fraud that failed at
the part this dissertation measures.

The strategy adopted is therefore \textbf{computational simulation with
analytical validation}: construct a formal model of the attack, compute
its consequences exactly where closed forms exist, simulate where they
do not, and validate the two against each other. The claims that result
are conditional --- ``under this model, with these parameters, the
following follows'' --- and the burden of the design is to make the
conditions explicit, the computation correct, and the results
reproducible.

Three commitments follow, and they shape everything below.

\textbf{Exactness before estimation.} Where a quantity has a closed
form it is computed exactly, not estimated. Attack success probability
is the binomial survival function
(\eqnref{eq:m5}), evaluated through the regularised incomplete beta
function; the coalition distribution is an exact convolution
(Section~\ref{sec:coalition-model}); the threshold solvers are exact
binary searches or bisections on monotone functions. Monte Carlo is
used to \emph{validate} these exact values and to reach quantities that
have no closed form, not as the primary method. This inverts the more
common simulation-study design, and it is what allows the Monte Carlo
error itself to be reported as a result
(Section~\ref{sec:findings-methods}).

\textbf{Falsifiability by construction.} Wherever a modelling decision
could flatter the conclusion, the design includes a mechanism that
would expose it. The Chernoff bound is checked against the exact value
in every trial, so a wrongly-implemented bound would show as a
violation rather than as a plausible number. The coalition convolution
is checked against a 200{,}000-draw Monte Carlo simulation in the test
suite. Reproducibility is enforced by a test that fails if two runs at
the same seed differ, and by a second test that fails if two runs at
\emph{different} seeds agree --- so the first test cannot pass
vacuously.

\textbf{Assumptions declared and swept.} Every modelling assumption is
recorded in \code{ASSUMPTIONS.md} and cross-referenced from the module
that relies on it. Where an assumption is both consequential and
ungrounded --- the relative costs of the three vectors, the capacity
ceilings --- it is not merely declared but swept, and the conclusion is
reported as a function of it (Section~\ref{sec:exp-sensitivity}).

\subsection{A two-layer architecture, and why both layers are necessary}
\label{sec:method-twolayer}

The implementation has two layers that answer the same questions by
different means.

The \textbf{concrete layer} is a working blockchain e-voting
application: SHA-256 block hashing, proof-of-work mining, chain
validation, voter registration, role-based election administration,
duplicate-vote prevention and result publication. Attacking it produces
an actual winner change on an actual mined chain.

The \textbf{abstract layer} is the mathematical framework: exact
binomial evaluation, Chernoff bounds, Monte Carlo, the threshold
solvers and the game-theoretic optimisers. It produces probabilities
rather than outcomes.

Neither layer alone would support the conclusions. The abstract layer
supplies statistical power --- \TotalTrials{} paired trials complete in
\RuntimeSeconds{} seconds, whereas mining that many elections is
infeasible --- but every one of its numbers is a probability computed
under a model, and a model can be wrong in ways that are invisible from
inside it. The concrete layer supplies ground truth --- a real tally, a
real winner, a real chain whose integrity is checked after every attack
--- but at a sample size (\ElectionTrials{} elections) too small for
the parameter sweeps.

Running both is a complementary analysis, not a validation of one
estimand by the other. The concrete experiment uses different electorate
and attack-resource ranges and derives its threshold from a realised
two-candidate margin, whereas E1 uses the abstract parameter distribution.
Section~\ref{sec:findings-concrete} therefore compares direction and
interpretation only; it does not claim that the two rates estimate the same
quantity.

\subsection{Variables and operationalisation}

\tabref{tab:variables} states how each construct is operationalised.
The distinction that matters most is between the three
\emph{independent} variables the experiment manipulates and the
\emph{control} variables that are held identical across the paired
arms.

\begin{table}[htbp]
\centering
\caption{Operationalisation of the model's variables.}
\label{tab:variables}
\small
\begin{tabular}{@{}p{2.5cm}p{2.6cm}p{8.4cm}@{}}
\toprule
\textbf{Role} & \textbf{Variable} & \textbf{Operationalisation} \\
\midrule
Independent & System configuration & Two levels: baseline (weak detection
and validation) and protected (blockchain-style audit, signatures,
identity verification, ballot validation --- modelled as higher
detection and lower acceptance probabilities). \\
\addlinespace
Independent & Attack vector & Three levels: bribery, Sybil
registration, ballot stuffing; plus the coalition of all three in
Section~\ref{sec:coalition}. \\
\addlinespace
Independent & Attacker resources & Budget $B$ and reward $R$, entering
through the feasibility constraint $T = \min(N, \lfloor B/R \rfloor)$. \\
\addlinespace
Dependent & Attack success probability & $\Prob(W \geq m)$, where $W$
counts fraudulent votes surviving every countermeasure. Computed
exactly; also estimated by Monte Carlo and bounded by Chernoff. \\
\addlinespace
Dependent & Election flip & Binary: did the winner of a concretely
mined election change? Measured only in the concrete layer. \\
\addlinespace
Dependent & Attacker expected utility & \eqnref{eq:utility}: payoff
weighted by flip probability, less cost, less expected penalty. \\
\addlinespace
Control & Scenario structure & $N, m, B, R, P, S, A$ --- identical
between the two arms of every pair by construction. \\
\addlinespace
Control & Random stream & A single seeded generator; identical seed
gives identical output bit-for-bit. \\
\bottomrule
\end{tabular}
\end{table}

\subsection{What the ledger contributes, and what it does not}
\label{sec:method-attribution}

Because the protected configuration bundles several controls together,
it is worth stating precisely which component does what before any
result is reported. The pipeline is ordered, and the ordering is the
whole point:

\begin{center}
\small
fraud attempt $\rightarrow$ \textbf{acceptance / identity / signature /
ballot validation} $\rightarrow$ \textbf{detection} $\rightarrow$
surviving votes only $\rightarrow$ \textbf{ledger records them}
\end{center}

Every fraudulent vote that is stopped in this model is stopped by the
validation or detection stage. The ledger sits at the end of the
pipeline and records whatever survives; it contributes integrity,
immutability, auditability and tamper evidence, and it contributes
nothing whatever to deciding \emph{whether a vote should have been
cast}. A blockchain cannot tell that a genuine voter was paid, that a
registration was forged, or that a ballot was injected --- those are
determined upstream, by mechanisms that would work equally well behind
a conventional append-only audit log.

This distinction is preserved in the experimental design rather than
merely asserted in the prose. The concrete layer runs the chain's own
validity check after every attack, so the two properties are measured
separately: the ledger's integrity, and the election's correctness. The
\emph{Findings} chapter reports that these come apart completely ---
integrity held in every trial, and elections were flipped anyway.

Nothing in this dissertation should therefore be read as showing that
``blockchain prevents bribery'', or Sybil attacks, or ballot stuffing.
The claim supported by the results is narrower and conditional: under
the assumptions of this model and over the sampled parameter ranges, a
protected architecture combining ledger-backed auditability with
identity verification, signature validation, ballot validation and
raised detection substantially reduces attack feasibility.

The dependent variable deserves one clarification. ``Attack success''
is defined throughout as reaching the manipulation threshold $m$, which
is the number of additional fraudulent votes needed to overturn the
margin (\eqnref{eq:threshold}). It is not defined as ``some fraud
occurred'' or ``fraud was undetected''. This is a deliberately
demanding definition: an attack that places a hundred fraudulent votes
into an election it fails to flip counts, here, as a failure.

\subsection{Parameter selection and its justification}
\label{sec:method-params}

Parameter ranges (\tabref{tab:parameters}) were selected to span
plausible deployment conditions rather than to be estimated from data,
because for most of these quantities no defensible estimate exists. The
probability that a forged voter registration is detected is not
published by any electoral authority in a form that would support a
point estimate. Three consequences follow, and the design accommodates
all three.

\textbf{Ranges, not points.} Every parameter is sampled from a range on
every trial, so no result depends on a single assumed value. The
electorate is drawn from the configured range 200--5{,}000, realising
\ElectorateMin{} to \ElectorateMax{} across the trials actually run, and the
threshold $m$ is itself randomised as 2--25\% of $N$, so the reported
means are averages over a distribution of scenarios rather than
properties of one configuration.

\textbf{Direction of bias stated.} Where a choice could flatter either
side it is made against the conclusion being argued. Legitimate ballots
always pass validation, so the model never credits the defence with
catching fraud by accident. Detection is applied independently per
item, which is more favourable to the attacker than correlated
detection, since an audit that catches one item usually catches its
batch.

\textbf{Sensitivity where it matters.} The coalition analysis depends
on four constants with no empirical basis --- two relative costs and
two capacity ceilings. Rather than defend particular values, the
analysis re-runs across \SensitivityConfigs{} combinations of them and
reports the conclusion's stability
(Section~\ref{sec:exp-sensitivity}).

The benchmark scenarios (S1--S5) are calibrated differently and for a
different purpose. For S1--S4 the threshold $m$ is placed close to the
baseline mean $\mu$ rather than at an extreme. This is necessary for
the three-method comparison to be informative: a threshold far below
$\mu$ makes success trivially certain, one far above makes it trivially
impossible, and the \emph{upper-tail} Chernoff form is not defined
when $m \leq \mu$ --- the gap that Section~\ref{sec:certificate}
closes with a lower-tail certificate. S5 is the deliberate exception --- a threshold far beyond what
its budget can reach --- and functions as a negative control: a model
that reported S5 as succeeding would be wrong in an obvious way.

\subsection{Paired randomised design}
\label{sec:method-paired}

The core experiment is a \textbf{matched-pairs design}. Each trial
draws one scenario and evaluates it twice, under both configurations,
with identical $N, m, B, R, P, S, A$ and identical random streams. Only
the acceptance and detection probabilities differ.

The alternative --- drawing independent samples for each arm --- would
confound the countermeasure effect with the difference between the
scenarios happening to be drawn for each arm. Because the per-trial
success probability has a standard deviation of \BriberyBaselineStd{}
for bribery, scenario-to-scenario variance dominates the effect being
measured, and an unpaired design would need far more trials to reach
the same precision. Pairing removes that variance entirely: the
reported reduction is a within-scenario effect, and every pair is its
own control.

\textbf{Inference.} Confidence intervals are computed by bootstrap
resampling of the paired differences (10{,}000 resamples) rather than
by a normal approximation. This is not a stylistic preference. The
per-trial success probabilities are strongly bimodal --- most scenarios
are either comfortably safe or comfortably lost --- which is visible in
the summary statistics: the baseline bribery distribution has a mean of
\BriberyBaseline{} but a median of \BriberyBaselineMedian. A
distribution with that shape violates the assumptions of a
$t$-interval, and the bootstrap makes no distributional assumption.

\textbf{Trial count.} \TotalTrials{} pairs per vector were run, above
the 1{,}000 originally specified, chosen so that the bootstrap
intervals on the mean reduction are narrow relative to the effect: the
bribery reduction interval $[\BriberyCILow, \BriberyCIHigh]$ has a
width under 0.04 against an effect of \BriberyAbsReduction.

\subsection{Triangulation of three computational methods}
\label{sec:method-triangulation}

Every trial computes its success probability three ways --- exact
binomial, Monte Carlo, and Chernoff bound where valid --- on identical
inputs. This is methodological triangulation, and each method plays a
distinct role.

\begin{description}[leftmargin=1.3em,style=nextline,itemsep=2pt]
\item[Exact binomial.] The ground truth. Every reported probability is
  this value; the other two exist to check it and to characterise their
  own usefulness.
\item[Monte Carlo.] An independent numerical path to the same quantity,
  through a different code path (sampling rather than the survival
  function). Agreement is evidence that neither contains an error;
  its error against the exact value quantifies how much simulation is
  needed for a given precision (RQ2).
\item[Chernoff bound.] A provable upper bound. Its role is not accuracy
  but \emph{certification}: a margin the bound rules out is ruled out
  regardless of simulation effort. Its looseness is measured and
  reported rather than treated as a defect.
\end{description}

The three are computed on the same trial, not on separate runs, which
is what makes the comparison a validation rather than a coincidence.
Critically, the bound is applied only where mathematically valid: the
dispatcher determines whether $m$ lies in the upper tail, the lower
tail, or neither, and returns no bound with a machine-readable reason
outside those regions. Averaging the three quantities across
heterogeneous trials --- as an earlier iteration of this work did ---
compares them on different sub-populations, since the bound exists only
where $m > \mu$; \figref{fig:detection} therefore holds every other
parameter fixed and sweeps only $q$.

\subsection{Threshold solver design}
\label{sec:method-thresholds}

The threshold solvers (T1--T7) invert the model, and three design
decisions govern their trustworthiness.

\textbf{Non-existence is a result.} Each solver returns a structured
object carrying either a value or an explicit absence with a
machine-readable reason. No solver substitutes an endpoint, a zero or a
sentinel for a threshold that does not exist, and the figures print
``does not exist'' rather than drawing a line. A threshold silently
defaulted to a domain boundary would be indistinguishable, in a plot,
from a real finding.

\textbf{Exactness where the domain is discrete.} T3, T6 and T7 search
over integers and use binary search, which is exact because each
underlying function is monotone in the search variable: $\Prob(X \geq
m)$ is non-decreasing in $n$, and the Chernoff bound is decreasing in
$m$ above the mean. These are not approximations.

\textbf{A definitional choice, stated.} T4 and T5 required resolving an
ambiguity in what ``the attacker's utility'' means. The headline
optimiser searches $T \geq 0$ and therefore includes declining to
attack, so its optimum is never negative --- and a quantity floored at
zero cannot be bisected for a zero crossing. The thresholds are
consequently defined on $\max_{T \geq 1} U(T)$, the signed payoff of
actually attacking, which goes negative exactly when attacking loses
money. This is a definition rather than a derivation, it changes the
reported numbers, and it is stated wherever those thresholds appear.

\subsection{Coalition search design}
\label{sec:method-coalition}

The coalition analysis searches the budget simplex exhaustively at a
0.05 step (231 allocations for three vectors) rather than using
continuous optimisation. The reason is structural: the objective is
piecewise flat wherever $\lfloor B/c \rfloor$ does not change, so
gradients are zero almost everywhere and undefined at the steps, which
defeats gradient methods. Exhaustive evaluation is affordable because
each evaluation is a convolution rather than a simulation. The reported
optimum is the best grid point and may differ slightly from the true
continuous optimum --- stated as a limitation rather than glossed.

The success probability of a coalition is computed exactly. A sum of
binomials with unequal success probabilities is not binomial and has no
closed form, and the obvious approximations (normal, Poisson) are least
accurate precisely in the tail that matters. Convolving the component
PMFs is exact, and truncating each at $m-1$ before convolving is also
exact rather than approximate, since $\Prob(W \geq m) = 1 - \Prob(W
\leq m-1)$ and no combination summing to at most $m-1$ can involve a
component larger than $m-1$. That truncation is what makes
national-scale margins tractable: it holds the transform length at $m$
rather than at the total attempt count, which is two orders of
magnitude larger.

\subsection{Validation and verification}
\label{sec:method-validation}

Correctness is established at four levels.

\textbf{Unit level (\TestCount{} tests).} Beyond routine coverage, several tests
encode properties the mathematics requires and that a plausible-looking
implementation can silently violate: that the lower-tail Chernoff bound
carries a negative exponent, since a positive one yields a ``bound''
exceeding 1 that is vacuous but looks reasonable in output; that the
bound is never returned outside its region of validity; that $T = 0$
rather than a division error when $B = 0$, $R = 0$ or $B < R$; and that
each threshold solver's answer actually brackets the crossing it claims
to find.

\textbf{Cross-method level.} The exact and Monte Carlo values are
compared on every trial; the Chernoff bound is checked against the
exact value on every trial where it is defined. A violation would
appear as a bound lying below the exact value it is supposed to bound
(\figref{fig:twosided}); none occurs across \BriberyChernoffValid{},
\SybilChernoffValid{} and \BallotStuffingChernoffValid{} valid trials
respectively.

\textbf{Cross-layer level.} The abstract probabilities and the concrete
mined-election flip rates are compared for ordering and magnitude
(Section~\ref{sec:findings-concrete}).

\textbf{Reproducibility level.} A test asserts that two runs at the same
seed produce identical output, and a second asserts that different
seeds do not, so the first cannot pass vacuously.

\subsection{Figure geometry on a discrete support}
\label{sec:method-figure-geometry}

Every random variable in this model is integer-valued: surviving
fraudulent votes, attempts, and margins are counts. That is a property
of the model with a consequence for how its distributions are drawn,
and it is recorded here because a figure can be numerically exact and
still mislead --- a failure mode no numerical test detects.

Plotting a discrete probability mass function as unit-width centred
bars places the bar for outcome $k$ on the interval $[k-0.5,\,k+0.5]$.
The event $\{X \geq m\}$ therefore occupies the region whose left edge
lies at $m - 0.5$, not at $m$. Drawing the labelled threshold line at
$m$ places it through the \emph{centre} of the first bar of the success
region, so the colour change appears to begin half a vote before the
threshold it is labelled with, even though the shaded probability mass
is exactly correct. An earlier draft of \figref{fig:pmf} contained
precisely this defect. It was found visually in review rather than by
any test, which is why the three tests that now pin it were written
afterwards.

The convention adopted throughout is that boundaries between integer
outcomes are drawn at half-integers and labelled with the integer
margin they separate: the manipulation threshold at $m - 0.5$, the
reachable margin $m_\alpha$ at $m_\alpha + 0.5$ (since margins at or
\emph{below} $m_\alpha$ are reachable), and the certified margin
$m_\varepsilon$ at $m_\varepsilon - 0.5$ (since margins at or
\emph{above} $m_\varepsilon$ are certified). Quantities that are
genuinely continuous are not moved: the mean $\mu = T(1-q)$ is an
expectation rather than a boundary between outcomes, and is drawn at
its true coordinate.

The correction changes no probability, threshold, table or finding in
this report: the shaded mass of \figref{fig:pmf} agreed with
$P(X \geq m)$ to machine precision both before the correction and
after it. Three tests in
\code{tests/test\_figure\_geometry.py} now pin the geometry by
interrogating the plotting library's rectangle and line coordinates
directly, so the visual property is verified in the same way as the
numerical ones rather than left to the eye.

\subsection{Reproducibility protocol}
\label{sec:reproducibility}

All randomness derives from a single seeded
\code{numpy.random.Generator}; no global random state is used anywhere,
and no unseeded standard-library random calls exist in the codebase.
Re-running with seed \Seed{} reproduces every number in this report
bit-for-bit --- verified by regenerating the full pipeline from a clean
results directory and comparing the summary output. The threshold and
coalition analyses are fully analytic and require no seed at all.

The report itself is part of the protocol. Every \emph{experimental
result} in the prose is injected at build time from
\code{results/json/} and \code{results/csv/} as a \LaTeX{} macro rather
than typed. A macro whose experiment has not been run causes the
document to fail compilation rather than to print a stale value, so the
text cannot silently disagree with the data it describes. The exceptions
are deliberate and are not results: configuration values quoted from
\code{config.py}, worked-example arithmetic shown at reduced precision
so that it reads on the page, and the analytic counter-examples of
Section~\ref{sec:tailforms}, which are computed and regression-tested in
\code{tests/test\_chernoff.py} rather than in the results pipeline.

\subsection{Ethical considerations}

This work models attacks on elections, which warrants explicit
treatment.

No real system, election, voter or personal datum is involved at any
point. All identifiers are simulation-internal indices; the candidates
are \emph{A} and \emph{B} and have no referent. The artefact is a
mathematical model, not a tool: it computes probabilities and contains
no capability that could be directed at a deployed system, and nothing
in it constitutes an exploit, a technique for evading a specific
control, or an implementation of one.

The disclosure balance favours publication. The vectors modelled ---
bribery, Sybil registration, ballot stuffing --- are long-established
in the literature reviewed in Section~\ref{sec:litreview}, and the
contribution is a quantification of their effect and of the
countermeasure strength required to blunt them. That quantification is
of substantially more use to a defender, who can act on a detection
threshold, than to an attacker, who learns only that the attacks they
already knew about work about as well as expected.
